\begin{proof}[Proof of \cref{obj:lem:bv-lipschitz}]
Fix \(u,v\in\mathbb R_+^4\), and write
\[
s_\eta(u)=\sum_{\omega=0}^1u_{\eta\omega},\qquad
s(u)=\sum_{\eta=0}^1s_\eta(u),\qquad
\Delta_{uv}=\sum_{\eta,\omega=0}^1|u_{\eta\omega}-v_{\eta\omega}|.
\]
In particular, \(|s(u)-s(v)|\leq\Delta_{uv}\). For either arm \(\eta\), the stabilized contribution in \cref{obj:def:dual-process} satisfies
\[
0\leq g_\eta(u)\leq s(u),\qquad
|g_\eta(u)-g_\eta(v)|
 \leq(1+\epsilon^{-1})\Delta_{uv}.
\]
Here is a derivation of the second bound, including zero masses. Put
\(m_u=s_\eta(u)\) and \(z_u=u_{\eta1}\), and define \(m_v,z_v\) analogously.
When \(s(u)>0\), the two branches are
\[
g_\eta(u)=
\begin{cases}
z_u/\epsilon,&m_u\leq\epsilon s(u),\\
s(u)z_u/m_u,&m_u\geq\epsilon s(u).
\end{cases}
\]
They agree at equality; when \(s(u)=0\), \(g_\eta(u)=0\). If both vectors use the first branch, their difference is at most
\(\epsilon^{-1}|z_u-z_v|\). If both use the second, \(m_u,m_v>0\), and
\[
\left|\frac{z_u}{m_u}-\frac{z_v}{m_v}\right|
\leq
\frac{|z_u-z_v|+
 |(m_u-z_u)-(m_v-z_v)|}{m_u}.
\]
Since \(z_v/m_v\leq1\) and \(s(u)/m_u\leq\epsilon^{-1}\), splitting the difference of the two products bounds it by
\(|s(u)-s(v)|+\epsilon^{-1}\Delta_{uv}\).

For the mixed case \(m_u\leq\epsilon s(u)\) and
\(m_v\geq\epsilon s(v)\), the identity
\[
\frac{z_u}{\epsilon}-s(v)\frac{z_v}{m_v}
=
\frac{z_u-z_v}{\epsilon}
+\frac{z_v}{m_v}\frac{m_v-\epsilon s(v)}{\epsilon}
\]
and the bounds
\[
0\leq m_v-\epsilon s(v)
\leq(m_v-m_u)-\epsilon\bigl(s(v)-s(u)\bigr),
\qquad 0\leq\frac{z_v}{m_v}\leq1
\]
give the upper bound below. The rearrangement
\[
(z_u-z_v)+\frac{z_v}{m_v}(m_v-m_u)
=\left(1-\frac{z_v}{m_v}\right)(z_u-z_v)
 +\frac{z_v}{m_v}\bigl((m_v-z_v)-(m_u-z_u)\bigr)
\]
shows that its left-hand side has absolute value at most \(\Delta_{uv}\), since \(0\leq z_v/m_v\leq1\). Consequently,
\[
\frac{z_u}{\epsilon}-s(v)\frac{z_v}{m_v}
\leq \frac{(z_u-z_v)+(z_v/m_v)(m_v-m_u)}{\epsilon}
 -\frac{z_v}{m_v}\bigl(s(v)-s(u)\bigr)
\leq(1+\epsilon^{-1})\Delta_{uv}.
\]
For the lower bound, return to the first mixed-case identity: its correction term \((z_v/m_v)(m_v-\epsilon s(v))/\epsilon\) is nonnegative. Hence the difference is at least \((z_u-z_v)/\epsilon\geq-\epsilon^{-1}\Delta_{uv}\). Interchanging \(u,v\) handles the other mixed case. Finally, if either total mass is zero, all its coordinates and its arm contribution vanish; \(0\leq g_\eta\leq s\) and \(|s(u)-s(v)|\leq\Delta_{uv}\) give \(|g_\eta(u)-g_\eta(v)|\leq(1+\epsilon^{-1})\Delta_{uv}\) directly. The same inequalities also establish \(0\leq g_\eta(u)\leq s(u)\) in every branch. % lean: CausalSmith.Stat.DiscreteBudgetvalueCurve.budget_armValue_lipschitz

Define the positive gain and its hinge by
\[
\gamma_u=[g_1(u)-g_0(u)]_+,\qquad
\Phi_u(\lambda)=[\gamma_u-\lambda s(u)]_+,
\]
and similarly for \(v\). The arm bounds give
\(0\leq\gamma_u\leq s(u)\) and \(0\leq\gamma_v\leq s(v)\).
Because \(\lambda s(u)\geq0\) on \([0,1]\),
\[
f_\lambda(u)=g_0(u)+\Phi_u(\lambda).
\]
Taking positive parts and using the arm contribution bound yields
\[
|\gamma_u-\gamma_v|
\leq |g_1(u)-g_1(v)|+|g_0(u)-g_0(v)|
\leq2(1+\epsilon^{-1})\Delta_{uv}.
\]
% lean: CausalSmith.Stat.DiscreteBudgetvalueCurve.budget_positiveGain_lipschitz

Set
\[
\rho_u=
\begin{cases}\gamma_u/s(u),&s(u)>0,\\0,&s(u)=0,\end{cases}
\qquad
\rho_v=
\begin{cases}\gamma_v/s(v),&s(v)>0,\\0,&s(v)=0,\end{cases}
\qquad
\Psi_{uv}(\lambda)=\Phi_u(\lambda)-\Phi_v(\lambda).
\]
Both breakpoints lie in \([0,1]\), including when a total mass vanishes. The positive-part inequality gives, throughout that interval,
\[
|\Psi_{uv}(\lambda)|
\leq |\gamma_u-\gamma_v|+|s(u)-s(v)|
=:\Xi_{uv}.
\]
If \(\rho_u\leq\rho_v\), the three successive affine expressions for
\(\Psi_{uv}\) are
\[
\gamma_u-\gamma_v+\lambda\bigl(s(v)-s(u)\bigr),\qquad
-\gamma_v+\lambda s(v),\qquad 0
\]
on \([0,\rho_u]\), \([\rho_u,\rho_v]\), and \([\rho_v,1]\), respectively.
If \(\rho_v\leq\rho_u\), they are
\[
\gamma_u-\gamma_v+\lambda\bigl(s(v)-s(u)\bigr),\qquad
\gamma_u-\lambda s(u),\qquad 0
\]
on \([0,\rho_v]\), \([\rho_v,\rho_u]\), and \([\rho_u,1]\).
The formulas agree at their shared endpoints, also for coincident breakpoints. Variation on each affine piece equals the absolute difference of its endpoint values, which is at most \(2\Xi_{uv}\) by the preceding pointwise bound. Adding the three pieces gives
\[
\operatorname{TV}_{[0,1]}(\Psi_{uv})
\leq6\bigl(|\gamma_u-\gamma_v|+|s(u)-s(v)|\bigr).
\]
% lean: CausalSmith.Stat.DiscreteBudgetvalueCurve.hinge_difference_variation_le

The constant \(g_0(u)-g_0(v)\) in
\(f_\lambda(u)-f_\lambda(v)=g_0(u)-g_0(v)+\Psi_{uv}(\lambda)\)
leaves variation unchanged. At \(\lambda=0\), the positive-part inequality and the arm contribution bound give
\[
|f_0(u)-f_0(v)|
\leq 2|g_0(u)-g_0(v)|+|g_1(u)-g_1(v)|
\leq\bigl(3(1+\epsilon^{-1})+1\bigr)\Delta_{uv}.
\]
Thus the positive-gain and variation bounds, together with \(|s(u)-s(v)|\leq\Delta_{uv}\), imply
\[
\|f_{\cdot}(u)-f_{\cdot}(v)\|_{BV([0,1])}
\leq\bigl(15(1+\epsilon^{-1})+7\bigr)\Delta_{uv}.
\]
The displayed finite, nonnegative coefficient supplies \(C_\epsilon\).
% lean: CausalSmith.Stat.DiscreteBudgetvalueCurve.thresholdFun_bv_lipschitz
\end{proof}
